% ===========================================================================
% Chapter 3: Information
% ===========================================================================
% Peter Gacs, "Descriptional Complexity and Randomness" (arXiv:2105.04704),
% Chapter 3, source lines 3757-4870.  His statements and proofs are reproduced
% verbatim, with the author's permission; the Lean annotations, the dependency
% edges, the formalization notes and the few proof environments marked "ours"
% are not his.  Lean side: Gacs_Chapter_3/ (see Gacs_Chapter_3/strategy.md).
%
% Numbering: as in content.tex, theorems run on one counter per section and all
% other numbered environments share a second one, so every theorem, lemma,
% definition and remark below carries the number printed in his manuscript.
% Equations are numbered per chapter here and per section in his manuscript.
%
% Labels: his \label{...} is kept.  A numbered item he did not label is
% labelled <kind>:g<number>, e.g. thm:g3.1.4 for his Theorem 3.1.4.
%
% References into Chapter 1 whose targets content.tex does not reproduce
% (Section 1.6, Subsection 1.6.3, equations (1.6.4), (1.6.13), (1.7.2)) are
% given as plain numbers.  References to his Chapter 2 keep his labels.

\chapter{Information}

This chapter reproduces Chapter~3 of G\'acs' manuscript~\cite{Gacs2021} and records how
each of its results is formalized in Lean. Every numbered definition, theorem, lemma,
proposition and corollary of the chapter is formalized, with no \texttt{sorry}, and so are
Remarks~\ref{rem:inf-def-rand} and~\ref{rem:g3.1.16}. Only informal asides are left out,
among them Remark~\ref{rem:g3.1.13} and the claim that \( K^{*} \) is not semicomputable;
the formalization notes say where. Theorem~\ref{thm:co-rarely-co} and
Lemma~\ref{lem:co-co-rare} are formalized in a sharper form than the printed one.

Some conventions of the formalization are used throughout. For an optimal prefix machine
\( U \) every complexity is finite, so the Lean code works with natural numbers:
\texttt{Kolmogorov.Information.Kc U x y} is \( K(x \mvert y) \) and
\texttt{Kolmogorov.Information.Ku U x} is \( K(x) \). A relation \( f \lea g \) is stated as
``there is a \( c \) with \( f \le g + c \) for all arguments'', over the integers when
differences occur (\texttt{Kolmogorov.Information.AddLE},
\texttt{Kolmogorov.Information.AddEq}). Pairs are encoded by \texttt{Kolmogorov.pairCode}, and
a number inside a complexity, such as \( K(x) \) in \( K(y \mvert x, K(x)) \), by the unary
code \texttt{Kolmogorov.natCode}; any other computable encoding changes the complexities
by \( O(1) \). Classical entropies take values in \( [0, \infty] \), so identities that
involve no subtraction hold with no finiteness hypothesis.

\section{Information-theoretic relations} 

In this subsection, we use some material
from~\cite{GacsTrompVitAlgorStat00}.

\subsection{The information-theoretic identity}

Information-theoretic entropy is related to complexity in both a formal and a
more substantial way.

\subsubsection{Classical entropy and classical coding}

Entropy is used in information theory in order to characterize the
compressibility of a statistical source.

% Gacs source line 3774
 \begin{definition}
  \label{def:g3.1.1}
  \lean{Kolmogorov.Information.Classical.entropy}
  \leanok
The \df{entropy} of a discrete probability distribution \( P \) is defined as
 \[
   \cH(P) = -\sum_{x} P(x)\log P(x).
 \]
 \end{definition}

\begin{formalization}
A distribution is a function \( P \colon \alpha \to [0, \infty] \) of total mass one on a
type \( \alpha \), and the entropy is measured in bits with values in \( [0, \infty] \), so a
distribution of infinite entropy is allowed.
\end{formalization}

A simple argument shows that this quantity is non-negative.
It is instructive to recall a more general inequality, taken
from~\cite{CsiszarKornerBook81}, implying it:

% Gacs source line 3784
 \begin{theorem}\label{thm:CsiszarConc}
  \lean{Kolmogorov.Information.Classical.log_sum_inequality, Kolmogorov.Information.Classical.log_sum_eq_iff, Kolmogorov.Information.Classical.log_sum_nonneg}
  \leanok
Let \( a_{i}, b_{i} > 0 \), \( a = \sum_{i} a_{i} \), \( b = \sum_{i} b_{i} \).
Then we have
 \begin{equation}\label{eq:CsiszarConc}
  \sum_{i} a_{i} \log \frac{a_{i}}{b_{i}} \ge a \log \frac{a}{b},
 \end{equation}
with equality only if \( a_{i}/b_{i} \) is constant.
In particular, if \( \sum_{i}a_{i}=1 \) and \( \sum_{i}b_{i}\le 1 \) then we have
\( \sum_{i} a_{i} \log \frac{a_{i}}{b_{i}}\ge 0 \).
 \end{theorem}

\begin{formalization}
Stated for a finite index set, with real logarithms. The equality case is proved as an
equivalence: equality holds if and only if \( a_{i}/b_{i} = a/b \) for all \( i \).
\end{formalization}

% Gacs source line 3795
\begin{proof}
  \proves{thm:CsiszarConc}
  \leanok
Exercise, an application of Jensen's inequality to the
concave function \( \log x \).
\end{proof}

\begin{formalization}
The Lean proof applies Jensen's inequality to the strictly convex function \( x \log x \)
(Mathlib), which also gives the equality case.
\end{formalization}

If \( P \) is a probability distribution and \( Q(x) \ge 0 \), 
\( \sum_{x} Q(x) \le 1 \) then this inequality implies
 \[
   K(P) \le \sum_{x} P(x) \log Q(x).
 \]
We can interpret this inequality as follows.
We have seen in Subsection~1.6.3 that if
\( x \mapsto z_{x} \) is a prefix code, mapping the objects \( x \) into binary
strings \( z_{x} \) that are not prefixes of each other then
\( \sum_{x} 2^{-\len{z_{x}}} \le 1 \).
Thus, substituting \( Q(x) = 2^{-\len{z_{x}}} \) above we obtain Shannon's result:

% Gacs source line 3811
 \begin{theorem}[Shannon]\label{thm:Shannon-coding-lb}
  \lean{Kolmogorov.Information.Classical.shannon_coding_lb}
  \leanok
  \uses{def:g3.1.1}
If \( x \mapsto z_{x} \) is a prefix code then for its expected codeword length
we have the lower bound:
 \[
   \sum_{x} P(x) \len{z_{x}} \ge K(P).
 \]
 \end{theorem}

\begin{formalization}
Here and in the inequality before the theorem, \( K(P) \) stands for the entropy
\( \cH(P) \), as in the next paragraph, and that inequality reads
\( \cH(P) \le -\sum_{x} P(x) \log Q(x) \) (Gibbs' inequality,
\texttt{Kolmogorov.Information.Classical.entropy\_le\_crossEntropy}). The Lean statement
bounds the expected length from below by \( \cH(P) \), for a code on a countable type in which
no codeword is a prefix of the codeword of a different object.
\end{formalization}

% Proof environment ours: Gacs derives the theorem in the paragraph above it.
\begin{proof}
  \proves{thm:Shannon-coding-lb}
  \leanok
  \uses{l:kraft}
Gibbs' inequality with \( Q(x) = 2^{-\len{z_{x}}} \), which satisfies
\( \sum_{x} Q(x) \le 1 \) by Kraft's inequality.
\end{proof}


\subsubsection{Entropy as expected complexity}
Applying Shannon's theorem~\ref{thm:Shannon-coding-lb}
to the code obtained by taking \( z_{x} \) as the
shortest self-delimiting description of \( x \), we obtain the inequality
 \[
   \sum_{x} P(x)K(x) \ge \cH(P)
 \]
On the other hand, since \( \m(x) \) is a universal semimeasure, we have
\( \m(x) \gem P(x) \): more precisely, \( K(x) \lea -\log P(x) + K(P) \), 
leading to the following theorem.

% Gacs source line 3831
 \begin{theorem}
  \label{thm:g3.1.3}
  \lean{Kolmogorov.Information.Classical.entropy_expectedComplexity, Kolmogorov.Information.Classical.entropy_le_expectedComplexity, Kolmogorov.Information.Classical.expectedComplexity_le}
  \leanok
  \uses{def:g3.1.1, def:K}
For a computable distribution \( P \) we have
 \begin{equation}\label{eq:exp-entr-compl}
   \cH(P) \le \sum_{x} P(x)K(x) \lea \cH(P) + K(P).
 \end{equation}
 \end{theorem}

\begin{formalization}
\( P \) is a computable probability distribution on strings
(\texttt{Kolmogorov.Randomness.ComputableMeasure}), and \( K(P) \) is the complexity
\( K(y_{P}) \) of any string \( y_{P} \) that describes \( P \) as a lower semicomputable
function (\texttt{Kolmogorov.Randomness.IsLSCDescription}). The constant does not depend on
\( P \).
\end{formalization}

% Proof environment ours: Gacs gives the argument in the paragraph above the theorem.
\begin{proof}
  \proves{thm:g3.1.3}
  \leanok
  \uses{thm:Shannon-coding-lb, l:kraft, thm:coding}
The lower bound is Gibbs' inequality for the lengths \( K(x) \), which satisfy Kraft's
inequality \( \sum_{x} 2^{-K(x)} \le 1 \). For the upper bound,
\( 2^{-K(P)} P(x) \lem 2^{-K(x)} \), so \( K(x) \le -\log P(x) + K(P) + c \) whenever
\( P(x) > 0 \); multiply by \( P(x) \) and sum.
\end{proof}

\begin{formalization}
The bound \( 2^{-K(P)} P(x) \lem 2^{-K(x)} \) is a weak form of
Theorem~\ref{thm:exact-domin} (which has \( \m(P) \gem 2^{-K(P)} \) in place of
\( 2^{-K(P)} \)). It is proved in the Chapter~2 files, uniformly in the description of
\( P \), as \texttt{Kolmogorov.Randomness.semimeasure\_le\_exact}.
\end{formalization}

Note that \( \cH(P) \) is the entropy of the distribution \( P \) while \( K(P) \) is
just the description complexity of the function \( x \mapsto P(x) \), it has
nothing to do directly
with the magnitudes of the values \( P(x) \) for each \( x \) as real numbers.
These two relations give
 \[ 
K(P) \eqa \sum_{x} P(x)K(x) - K(P),
 \]
the entropy is within an additive constant equal to the
\df{expected complexity}.
Our intended interpretation of \( K(x) \) as information content of the
individual object \( x \) is thus supported by a tight quantitative
relationship to Shannon's statistical concept.


\subsubsection{Identities, inequalities}
The statistical entropy obeys meaningful identities which immediately
follow from the definition of conditional entropy. 

% Gacs source line 3856
 \begin{definition}
  \label{def:g3.1.2}
  \lean{Kolmogorov.Information.Classical.condEntropy}
  \leanok
  \uses{def:g3.1.1}
Let \( X \) and \( Y \) be two discrete random variables with a joint distribution.
 The \emph{conditional entropy} of \( Y \) with respect to \( X \) is defined as
 \[
     \cH(Y \mvert X) =
   -\sum_{x} P[X=x]\sum_y P[Y=y \mvert X=x]\log P[Y=y \mvert X=x].
 \]
The \df{joint entropy} of \( X \) and \( Y \) is defined as the entropy of
the pair \( (X,Y) \).
 \end{definition}

\begin{formalization}
The pair \( (X, Y) \) is given by its joint law \( p \colon \alpha \times \beta \to [0, \infty] \),
and the laws of \( X \) and \( Y \) are its marginals
(\texttt{Kolmogorov.Information.Classical.marg1}, \texttt{Kolmogorov.Information.Classical.marg2}).
The joint entropy is the entropy of \( p \).
\end{formalization}

The following \df{additivity property} is then verifiable by simple calculation:
 \begin{equation}\label{eq:entraddit}
   \cH(X,Y) = \cH(X) + \cH(Y \mvert X).
 \end{equation}
  Its meaning is that the amount of information in the pair \( (X,Y) \) is
equal to the amount of information in \( X \) plus the amount of our
residual uncertainty about \( Y \) once the value of \( X \) is known.

Though intuitively expected, the following identity needs proof:
 \begin{align*}
   \cH(Y\mvert X) \le \cH(Y).
 \end{align*}
We will prove it after we define the difference below.

\begin{formalization}
Additivity \eqref{eq:entraddit} is
\texttt{Kolmogorov.Information.Classical.entropy\_eq\_add\_condEntropy}, an identity in
\( [0, \infty] \) with no finiteness hypothesis. The inequality
\( \cH(Y \mvert X) \le \cH(Y) \) is
\texttt{Kolmogorov.Information.Classical.condEntropy\_le\_entropy}, proved directly from
Gibbs' inequality, fibre by fibre.
\end{formalization}

% Gacs source line 3881
 \begin{definition}
  \label{def:g3.1.3}
  \lean{Kolmogorov.Information.Classical.mutualInfo}
  \leanok
  \uses{def:g3.1.1, def:g3.1.2}
The \df{information} in \( X \) about \( Y \) is defined by
 \[
 \cI(X : Y) = \cH(Y) - \cH(Y \mvert X).
 \]
\end{definition}

\begin{formalization}
The difference is taken in \( [0, \infty] \), so it is meaningful when
\( \cH(Y) < \infty \), the case in which the identities below are stated.
\end{formalization}

This is the amount by which our knowledge of \( X \) decreases our 
uncertainty about \( Y \).

The identity below comes from equation~\eqref{eq:entraddit}.
 \[
 \cI(X : Y) = \cI(Y : X) = \cH(X) + \cH(Y) - \cH(X,Y)
 \]
The meaning of the symmetry relation is 
that the \emph{quantity} of information in \( Y \) about
\( X \) is equal to the quantity of information in \( X \) about \( Y \). 
Despite its simple derivation, this fact is not very intuitive; especially
since only the quantities are equal, the actual contents are generally
not (see~\cite{GacsKornComnInf73}). 
We also call \( \cI(X : Y) \) the \df{mutual information} of \( X \) and \( Y \).
We can also write
 \begin{align}\label{eq:info-via-exp-val}
   \cI(X : Y) = \sum_{x,y}P[X=x,Y=y]\log\frac{P[X=x,Y=y]}{P[X=x]P[Y=y]}.
 \end{align}

\begin{formalization}
Formalized when \( \cH(X), \cH(Y) < \infty \): the symmetry
\texttt{Kolmogorov.Information.Classical.mutualInfo\_swap}, the identity
\( \cI(X : Y) = \cH(X) + \cH(Y) - \cH(X,Y) \)
(\texttt{Kolmogorov.Information.Classical.mutualInfo\_eq\_entropy\_add\_sub}), and
\eqref{eq:info-via-exp-val}
(\texttt{Kolmogorov.Information.Classical.mutualInfo\_eq\_mutualInfoKL}). The right-hand side
of \eqref{eq:info-via-exp-val} is written as the sum of the non-negative terms
\( p_{1}(x) p_{2}(y)\, \phi\bigl(p(x,y)/(p_{1}(x) p_{2}(y))\bigr) / \ln 2 \) with
\( \phi(t) = t \ln t + 1 - t \) (Mathlib's \texttt{InformationTheory.klFun}), which has the
same sum (\texttt{Kolmogorov.Information.Classical.mutualInfoKL}).
\end{formalization}

% Gacs source line 3906
\begin{proposition}
  \label{prop:g3.1.4}
  \lean{Kolmogorov.Information.Classical.mutualInfo_eq_zero_iff}
  \leanok
  \uses{def:g3.1.3}
  The information \( \cI(X:Y) \) is nonnegative and is 0 only if \( X,Y \) are independent.
\end{proposition}

\begin{formalization}
Non-negativity is automatic for a value in \( [0, \infty] \). The Lean statement is an
equivalence, under \( \cH(X), \cH(Y) < \infty \): \( \cI(X : Y) = 0 \) if and only if
\( p(x,y) = p_{1}(x) p_{2}(y) \) for all \( x, y \).
\end{formalization}

% Gacs source line 3909
\begin{proof}
  \proves{prop:g3.1.4}
  \leanok
  \uses{thm:CsiszarConc}
Apply Theorem~\ref{thm:CsiszarConc} to~\eqref{eq:info-via-exp-val}.
This is 0 only if \( P[X=x,Y=y]=P[X=x]P[Y=y] \) for all \( x,y \).
\end{proof}

\begin{formalization}
The Lean proof uses the non-negative form of \eqref{eq:info-via-exp-val} instead of
Theorem~\ref{thm:CsiszarConc}: the sum vanishes if and only if every term does, and
\( \phi(t) = 0 \) only for \( t = 1 \).
\end{formalization}


\subsubsection{The algorithmic addition theorem}
For the notion of algorithmic entropy \( K(x) \) defined in
Section~1.6, it is a serious test of fitness to ask whether it
has an additivity property analogous to the identity~\eqref{eq:entraddit}. 
Such an identity would express some deeper
relationship than the trivial identity~\eqref{eq:entraddit} since the
conditional complexity is defined using the notion of 
\emph{conditional computation}, and not by an algebraic formula involving
the unconditional definition.

A literal translation of the identity~\eqref{eq:entraddit} turns out to
be true for the function \( K(x) \) as well as for Kolmogorov's
complexity \( C(x) \) with good approximation. 
But the exact additivity property of algorithmic entropy is subtler. 

% Gacs source line 3930
\begin{theorem}[Addition]
  \label{thm:g3.1.4}
  \lean{Kolmogorov.Information.addition, Kolmogorov.KPPair_chain_upper, Kolmogorov.KPPair_chain_lower}
  \leanok
  \uses{def:K, def:pair-complexity}
We have
 \begin{equation}\label{eq:Haddit}
	K(x, y) \eqa K(x) + K(y \mvert x, K(x)).
 \end{equation}
 \end{theorem}

\begin{formalization}
The condition \( x, K(x) \) is the string \texttt{pairCode x (natCode k)} with
\( k = K(x) \) (\texttt{Kolmogorov.Information.xk}). The two directions are proved upstream, for the
complexity \( K(x, y) \) of \texttt{Kolmogorov.pairCode x y}
(\texttt{Kolmogorov.KPPair\_chain\_upper}, \texttt{Kolmogorov.KPPair\_chain\_lower}).
\texttt{Kolmogorov.Information.addition} restates them with natural-number complexities.
\end{formalization}

% Gacs source line 3935
 \begin{proof}
  \proves{thm:g3.1.4}
  \leanok
  \uses{l:projmeas, thm:coding}
  We first prove \( \lea \).
Let \( F \) be the following s.d.~interpreter.
The machine computing \( F(p) \) tries to parse \( p \) into \( p = uv \) so that
\( T(u) \) is defined, then outputs the pair 
\( \tup{ T(u), T(v,\ang{ T(u), \len{u} } )} \). 
If \( u \) is a shortest description of \( x \) and \( v \) a
shortest description of \( y \) given the pair \( \tup{ x,K(x)} \) then the
output of \( F \) is the pair \( \tup{ x,y}  \). 

Now we prove \( \gea \). 
Let \( c \) be a constant (existing in force of 
Lemma~\ref{l:projmeas}) such that for all \( x \) we have 
\( \sum_y2^{-K(x,y)} \le 2^{-K(x)+c} \). 
Let us define the semicomputable function \( f(z,y) \) by 
 \[	
  f(\tup{ x,k} ,y) = 2^{k - K(x,y) - c}.
 \]
  Then for \( k=K(x) \) we have \( \sum_y f(z,y) \le 1 \), hence the
generalization~(1.6.13) of the coding theorem implies
 \(  K(y \mvert x, K(x)) \lea K(x,y) - K(x)  \).
 \end{proof}

Recall the definition of \( x^{*} \) as the first shortest description of \( x \).

\begin{formalization}
\( x^{*} \) is \texttt{Kolmogorov.Information.xStar}, \emph{some} shortest description of
\( x \), chosen without any computability claim. Everything below uses \( x^{*} \) only
through \( x^{*} \asymp \ang{x, K(x)} \), proved after Definition~\ref{def:g3.1.6}, and that
holds for every choice.
\end{formalization}

% Gacs source line 3960
 \begin{corollary}
  \label{cor:g3.1.5}
  \lean{Kolmogorov.Information.Ku_xStar_pair}
  \leanok
  \uses{def:K, def:pair-complexity}
 We have
 \begin{equation}\label{eq:star-joint}
  K(x^{*}, y) \eqa K(x, y).
 \end{equation}
 \end{corollary}

% Proof environment ours: Gacs states the corollary without proof.
\begin{proof}
  \proves{cor:g3.1.5}
  \leanok
  \uses{thm:g3.1.4, cor:shortest-descriptions}
The pair \( \ang{x, y} \) is computable from \( \ang{x^{*}, y} \), since \( x = T(x^{*}) \);
this gives \( \gea \). Conversely,
\( K(x^{*}, y) \lea K(x^{*}) + K(y \mvert x^{*}) \) and \( K(x^{*}) \eqa K(x) \), so
\( K(x^{*}, y) \lea K(x) + K(y \mvert x^{*}) \eqa K(x, y) \) by~\eqref{eq:Haddit.alt} below.
\end{proof}

The Addition Theorem implies 
 \[ 	K(x,y) \lea K(x) + K(y \mvert x) \lea K(x) + K(y) 
 \]
  while these relations do not hold for Kolmogorov's complexity
\( C(x) \), as pointed out in the discussion of Theorem~\ref{thm:addit}. 
For an ordinary interpreter, extra information is needed to separate
the descriptions of \( x \) and \( y \).  
The self-delimiting interpreter
accepts only a description that provides the neccesary information to
determine its end. 

The term \( K(x) \) cannot be omitted from the condition in \( K(y \mvert x, K(x)) \) 
in the identity~\eqref{eq:Haddit}. 
Indeed, it follows from
equation~(1.7.2) that with \( y=K(x) \), we have \( K(x,y)-K(x) \eqa 0 \).
But \( K(K(x) \mvert x) \) can be quite large, as shown in the following
theorem given here without proof.

% Gacs source line 3985
\begin{theorem}[see~\protect\cite{GacsSymm74}]\label{thm:compl-compl}
  \lean{Kolmogorov.Information.CompComp.exists_KK_ge}
  \leanok
  \uses{def:K}
 There is a constant \( c \) such that for all \( n \)
there is a binary string \( x \) of length \( n \) with
 \[  K(K(x) \mvert x) \ge \log n - \log\log n -c.
 \]
 \end{theorem}

\begin{formalization}
Logarithms are \( \lfloor \log_{2} \cdot \rfloor \) (\texttt{Nat.log 2}), and \( K(x) \) is
written in the unary code \texttt{Kolmogorov.natCode} inside the complexity. The proof is the
one of Section~\ref{sec:compl-compl}.
\end{formalization}

We will prove this theorem in Section~\ref{sec:compl-compl}.
For the study of information relations, let us introduce yet another
notation.

% Gacs source line 3995
 \begin{definition}
  \label{def:g3.1.6}
  \lean{Kolmogorov.Information.InfoLE, Kolmogorov.Information.InfoLEMod, Kolmogorov.Information.InfoEquiv}
  \leanok
  \uses{def:K}
For any functions \( f,g,h \) with values in \( \dS \), let
 \[  f \preceq g \bmod h
 \]
  mean that there is a constant \( c \) such that \( K(f(x) \mvert g(x),h(x))
\le c \) for all \( x \).
We will omit \( \bmod h \) if \( h \) is not in the condition.
The sign \( \asymp \) will mean inequality in both directions.
 \end{definition}

\begin{formalization}
The functions are defined on an arbitrary index type. That \( f \asymp g \) allows replacing
\( f \) by \( g \) inside a complexity, as condition or as object, is
\texttt{Kolmogorov.Information.InfoEquiv.cond} and \texttt{Kolmogorov.Information.InfoEquiv.obj}.
\end{formalization}

The meaning of \( x \asymp y \) is that the objects \( x \) and \( y \) hold
essentially the same information. 
Theorem~\ref{thm:rectransf}
implies that if \( f\asymp g \) then we can replace \( f \) by \( g \) wherever
it occurs as an argument of the functions \( H \) or \( K \).  
As an illustration of this definition, let us prove the relation
 \[ 	x^{*} \asymp \ang{ x, K(x) } 
 \]
  which implies \( K(y \mvert x, K(x)) \eqa K(y \mvert x^{*}) \), and hence 
 \begin{equation} 	K(x,y) - K(x) \eqa K(y \mvert x^{*}).
\label{eq:Haddit.alt} 
 \end{equation}
  The direction \( \preceq \) holds because we can compute \( x=T(x^{*}) \)
and \( K(x)=\len{x^{*}} \) from \( x^{*} \). 
On the other hand, if we have \( x \)
and \( K(x) \) then we can enumerate the set \( E(x) \) of all shortest
descriptions of \( x \). 
By the corollary of Theorem~\ref{thm:Hstat}, the
number of elements of this set is bounded by a constant. 
The estimate~(1.6.4) implies \( K(x^{*} \mvert x,K(x)) \eqa 0 \).

\begin{formalization}
\( x^{*} \asymp \ang{x, K(x)} \) is \texttt{Kolmogorov.Information.xStar\_equiv}, and
\eqref{eq:Haddit.alt} is \texttt{Kolmogorov.Information.addition\_star}. The bound
\( K(x^{*} \mvert x, K(x)) \eqa 0 \) is obtained from
Corollary~\ref{cor:shortest-descriptions} through the conditional coding theorem
(\texttt{Kolmogorov.Information.condCoding\_trim}), applied to the uniform distribution on the
shortest descriptions of \( x \).
\end{formalization}

Though the string \( x^{*} \) and the pair \( \tup{ x, K(x) } \) contain the
same information, they are not equivalent under some stronger
criteria on the decoder. 
To obtain \( x \) from the shortest description
may take extremely long time. 
But even if this time \( t \) is
acceptable, to find any description of length \( K(x) \) for \( x \) might
require about \( t2^{K(x)} \) steps of search. 

The Addition Theorem can be made formally analogous to its 
information-theoretic couterpart using Chaitin's notational trick. 

% Gacs source line 4038
 \begin{definition}
  \label{def:g3.1.7}
  \lean{Kolmogorov.Information.Kstar}
  \leanok
  \uses{def:K}
Let \( K^{*}(y \mvert x) = K(y \mvert x^{*}) \).
 \end{definition}

Of course, \( K^{*}(x) = K(x) \).
Now we can formulate the relation~\eqref{eq:Haddit.alt} as
 \[	K^{*}(x, y) \eqa K^{*}(x) + K^{*}(y \mvert x).
 \]
  However, we prefer to use the function \( K(y \mvert x) \) since the
function \( K^{*}(y \mvert x) \) is not semicomputable.
Indeed, if there was
a function \( h(x,y) \) upper semicomputable in \( \ang{ x,y} \) such that
\( K(x,y) - K(x) \eqa h(x,y) \) then by the argument of the proof of the
Addition Theorem, we would have \( K(y \mvert x) \lea h(x,y) \) which is not
true, according to remark after the proof of the Addition Theorem.

\begin{formalization}
The identity \( K^{*}(x, y) \eqa K^{*}(x) + K^{*}(y \mvert x) \) is
\texttt{Kolmogorov.Information.Kstar\_addition}. The claim that \( K^{*} \) is not
semicomputable is not formalized.
\end{formalization}

% Gacs source line 4054
\begin{definition}
  \label{def:g3.1.8}
  \lean{Kolmogorov.Information.Info}
  \leanok
  \uses{def:K}
We define the algorithmic mutual information with the same formula as classical
information:
 \[
   I(x : y) = K(y) - K(y \mvert x).
 \]
\end{definition}

Since the quantity \(  I(x : y) \) is not quite equal to
\( K(x) - K(x \mvert y) \) (in fact, Levin showed (see~\cite{GacsSymm74}) that
they are not even asymptotically equal), we prefer to define information also by a
symmetrical formula.

% Gacs source line 4067
 \begin{definition}\label{def:I*}
  \lean{Kolmogorov.Information.InfoStar}
  \leanok
  \uses{def:K, def:pair-complexity}
The mutual information of two objects \( x \) and \( y \) is
 \[
	I^{*}(x : y) = K(x) + K(y) - K(x, y).
 \]
 \end{definition}

\begin{formalization}
Both informations are integer valued.
\end{formalization}

The Addition Theorem implies the identity
 \[
 	I^{*}(x : y) \eqa I(x^{*} : y) \eqa I(y^{*} : x).
 \]

\begin{formalization}
\texttt{Kolmogorov.Information.infoStar\_eq\_info\_xStar} and
\texttt{Kolmogorov.Information.infoStar\_eq\_info\_yStar}.
\end{formalization}

Let us show that classical information is indeed the expected value of
algorithmic information, for a computable probability distribution \( p \):

% Gacs source line 4081
\begin{lemma}
  \label{lem:g3.1.10}
  \lean{Kolmogorov.Information.expected_infoStar}
  \leanok
  \uses{def:g3.1.3, def:I*}
Given a computable joint probability mass distribution \( p(x,y) \) over \( (x,y) \)
we have 
\begin{align}\label{eq:eqamipmi}
\cI(X : Y) - K(p) & \lea  \sum_{x} \sum_y p(x,y) I^{*}(x : y) 
\\& \lea \cI(X : Y) + 2 K(p) , 
\nonumber
\end{align}
where \( K(p) \) is the length of the shortest prefix-free program that computes
\( p(x,y) \) from input \( (x,y) \).
\end{lemma}

\begin{formalization}
\( p \) is a probability distribution on pairs of strings, and \( K(p) \) is the complexity
\( K(y_{p}) \) of any string \( y_{p} \) describing \( \ang{x, y} \mapsto p(x, y) \) as a lower
semicomputable function (\texttt{Kolmogorov.Randomness.IsLSCDescription}). \( \cI(X : Y) \)
is the non-negative form \eqref{eq:info-via-exp-val}
(\texttt{Kolmogorov.Information.Classical.mutualInfoKL}). Both inequalities are stated
without subtraction, in \( [0, \infty] \), with
\( \sum p\, I^{*} = \sum p(x,y) K(x) + \sum p(x,y) K(y) - \sum p(x,y) K(x,y) \) rearranged,
so no finiteness hypothesis is needed.
\end{formalization}

% Gacs source line 4092
\begin{proof}
  \proves{lem:g3.1.10}
  \leanok
  \uses{thm:g3.1.3}
We have
\begin{align*}
\sum_{x,y} p(x,y) I^{*}(x : y) &\eqa 
\sum_{x,y} p(x, y)  [K(x) + K(y) - K(x, y)].
\end{align*}
 Define \( \sum_y p(x,y) = p_{1}(x) \) and \( \sum_{x} p(x,y) = p_{2}(y) \)
to obtain
\begin{align*}
\sum_{x,y} p(x,y) I(x:y) &\eqa
 \sum_{x} p_{1} (x) K(x) + \sum_y p_{2} (y) K(y) - \sum_{x,y} p(x,y) K(x, y).
\end{align*}
The distributions \( p_{i} \) (\( i=1,2 \)) are computable.
We have seen in~\eqref{eq:exp-entr-compl} that
\( \cH(q) \lea \sum_{x} q(x) K(x) \lea \cH(q) + K(q) \).

Hence, \( \cH(p_{i}) \lea \sum_{x} p_{i} (x) K(x) \lea \cH(p_{i}) + K(p_{i}) \).
(\( i=1,2 \)), and \( \cH(p) \lea \sum_{x,y} p(x,y) K(x,y) \lea \cH(p) + K(p) \).
On the other hand, the probabilistic mutual information
is expressed in the entropies by \( \cI(X : Y) = \cH(p_{1}) + \cH(p_{2}) - \cH(p) \).
By construction of the \( q_{i} \)'s above,
we have \( K(p_{1}), K(p_{2}) \lea K(p) \). 
Since the complexities are positive, substitution establishes the lemma.
\end{proof}

\begin{formalization}
The Lean proof avoids descriptions of the marginals and the bounds
\( K(p_{1}), K(p_{2}) \lea K(p) \). The joint law and its two marginals form jointly lower
semicomputable families indexed by the description \( y_{p} \), so indexed domination
(\texttt{Kolmogorov.Randomness.indexed\_domination\_trim}) gives
\( 2^{-K(y_{p})} p_{i}(x) \lem 2^{-K(x)} \) directly. The expected complexity bounds then
follow as in Theorem~\ref{thm:g3.1.3}.
\end{formalization}

% Gacs source line 4117
 \begin{remark}\label{rem:inf-def-rand}
The information \( I^{*}(x : y) \) can be written as
 \[
   I^{*}(x : y) \eqa \log \frac{\m(x, y)}{\m(x)\m(y)}.
 \]
Formally, \( I^{*}(x : y) \) looks like 
\( \ol\d_{\m\times\m}((x, y) \mvert \m \times \m) \) with the function 
\( \ol\d_{\cdot}(\cdot) \) introduced in~\eqref{eq:d_c}.
Thus, it looks like \( I^{*}(x : y) \) measures the deficiency of randomness
with respect to the distribution \( \m \times \m \).
The latter distribution expresses our ``hypothesis'' that \( x,y \) are
``independent'' from each other.
There is a serious technical problem with this interpretation: the function
\( \ol\d_{P}(x) \) was only defined for computable measures \( P \).
Though the definition can be extended, it is not clear at all that the
expression \( \ol\d_{P}(x) \) will also play the role of universal test
for arbitrary non-computable distributions.
Levin's theory culminating in~\cite{LevinRandCons84} develops this idea,
and we return to it in later parts of these notes.
 \end{remark}

\begin{formalization}
The displayed identity is formalized multiplicatively, for any universal semimeasure
\( \m \): \( 2^{-c} \m(x, y) \le 2^{I^{*}(x : y)} \m(x) \m(y) \le 2^{c} \m(x, y) \)
(\texttt{Kolmogorov.Information.infoStar\_eq\_log\_ratio}). The interpretation as a
deficiency of randomness is not formalized.
\end{formalization}

Let us examine the size of the
defects of naive additivity and information symmetry.

% Gacs source line 4149
 \begin{corollary}
  \label{cor:g3.1.12}
  \lean{Kolmogorov.Information.additivity_defect, Kolmogorov.Information.additivity_defect_le, Kolmogorov.Information.info_asymmetry}
  \leanok
  \uses{def:g3.1.8}
For the defect of additivity, we have
 \begin{equation}
  K(x) + K(y \mvert x) - K(x, y) \eqa I(K(x) : y \mvert x)
 \lea K(K(x) \mvert x).
  \end{equation}
For information asymmetry, we have
 \begin{equation}
 I(x : y) - I(y : x) \eqa I(K(y) : x \mvert y) - I(K(x) : y \mvert x).
 \end{equation}
 \end{corollary}

\begin{formalization}
The conditional information is
\( I(u : v \mvert w) = K(v \mvert w) - K(v \mvert w, u) \)
(\texttt{Kolmogorov.Information.CondInfo}).
\end{formalization}

% Gacs source line 4159
 \begin{proof}
  \proves{cor:g3.1.12}
  \leanok
  \uses{thm:g3.1.4}
Immediate from the addition theorem.
 \end{proof}

% Gacs source line 4163
 \begin{theorem}[Levin]
  \label{thm:g3.1.6}
  \lean{Kolmogorov.Information.levin_asymmetry, Kolmogorov.Information.asym_gap}
  \leanok
  \uses{def:g3.1.8}
The information \( I(x : y) \) is not even asymptotically symmetric.
 \end{theorem}

\begin{formalization}
Made precise as: there is a \( c \) such that for every \( n \) some \( x \) of length \( n \)
has, for one of the pairs \( (u, v) = (x, K(x)) \) or \( (x^{*}, K(x)) \),
\( 2\,|I(u : v) - I(v : u)| \ge \log n - \log\log n - c \).
\end{formalization}

% Gacs source line 4166
 \begin{proof}
  \proves{thm:g3.1.6}
  \leanok
  \uses{thm:compl-compl, thm:g3.1.4}
For some constant \( c \), assume
\( |x| = n \) and \( K(K(x) \mvert x) \ge \log n - \log\log n - c \).
Consider \( x, x^{*} \) and \( K(x) \).
Then we can check immediately the relations
 \[
  I(x : K(x)) \lea 3 \log \log n 
              < \log n - \log\log n - c \lea I(x^{*} : K(x)).
 \]
and
 \[
    I(K(x) : x) \eqa I(K(x) : x^{*}).
 \]
It follows that symmetry breaks down in an exponentially great measure
either on the pair \( (x, K(x)) \) or the pair \( (x^{*}, K(x)) \).
 \end{proof}

\begin{formalization}
The Lean proof does not bound \( I(x : K(x)) \) by \( 3 \log\log n \); it compares the two
asymmetries directly. Since \( K(K(x) \mvert x^{*}) \eqa 0 \) and
\( I(K(x) : x) \eqa I(K(x) : x^{*}) \), the asymmetry of \( (x^{*}, K(x)) \) minus that of
\( (x, K(x)) \) is \( \gea K(K(x) \mvert x) \ge \log n - \log\log n - c \)
(\texttt{Kolmogorov.Information.asym\_gap}), so one of the two is at least half of that.
\end{formalization}


\subsubsection{Data processing}
The following useful identity is also classical, and is called the
\df{data processing identity}:
 \begin{equation}\label{eq:data-proc-ident}
   \cI(Z : (X, Y)) = \cI(Z : X) + \cI(Z : Y \mvert X).
 \end{equation}
Let us see what corresponds to this for algorithmic information.

\begin{formalization}
The classical identity \eqref{eq:data-proc-ident} is
\texttt{Kolmogorov.Information.Classical.mutualInfo\_chain}, for a joint law of three
variables with \( \cH(X), \cH(Y) < \infty \), the conditional information being
\texttt{Kolmogorov.Information.Classical.condMutualInfo}.
\end{formalization}

Here is the correct conditional version of the addition theorem:
% Gacs source line 4193
 \begin{theorem}
  \label{thm:g3.1.7}
  \lean{Kolmogorov.Information.condAddition, Kolmogorov.KPCondPair_chain_upper, Kolmogorov.KPCondPair_chain_lower}
  \leanok
  \uses{def:K, def:pair-complexity}
We have
 \begin{equation}\label{eq:cond-addit}
  K(x, y \mvert z) \eqa K(x \mvert z) + K(y \mvert x, K(x \mvert z), z).
 \end{equation}
 \end{theorem}

\begin{formalization}
The condition of the last term is the string \( \ang{z, x, K(x \mvert z)} \)
(\texttt{Kolmogorov.Information.condCtx}); the order of the components changes complexities
only by \( O(1) \). The two directions are proved upstream
(\texttt{Kolmogorov.KPCondPair\_chain\_upper}, \texttt{Kolmogorov.KPCondPair\_chain\_lower}).
\end{formalization}

% Gacs source line 4199
 \begin{proof}
  \proves{thm:g3.1.7}
  \leanok
  \uses{l:projmeas, thm:coding}
The same as for the unconditional version.
 \end{proof}

% Gacs source line 4203
 \begin{remark}
  \label{rem:g3.1.13}
The reader may have guessed the inequality
 \[
  K(x, y \mvert z) \eqa K(x \mvert z) + K(y \mvert x, K(x), z),
 \]
but it is incorrect: taking \( z = x \), \( y = K(x) \),
the left-hand side equals \( K(x^{*} \mvert x) \), and the right-hand side
equals \( K(x \mvert x) + K(K(x) \mvert x^{*}, x) \eqa 0 \).
 \end{remark}

\begin{formalization}
Not formalized.
\end{formalization}

The following inequality is a useful tool.
It shows that conditional complexity is a kind of one-sided distance, and
it is the algorithmic analogue of the well-known (and not difficult
classical inequality)
 \[
   \cH(X \mvert Y) \le \cH(Z \mvert Y) + \cH(X \mvert Z).
 \]

\begin{formalization}
The classical inequality is
\texttt{Kolmogorov.Information.Classical.condEntropy\_triangle}, in \( [0, \infty] \).
\end{formalization}

% Gacs source line 4221
 \begin{theorem}[Simple triangle inequality]\label{thm:simple-triangle}
  \lean{Kolmogorov.Information.simple_triangle}
  \leanok
  \uses{def:K, def:pair-complexity}
We have
 \begin{equation}\label{eq:simple-triangle}
  K(z \mvert x) \lea K(y, z \mvert x) \lea K(y \mvert x) + K(z \mvert y).
 \end{equation}
 \end{theorem}

% Gacs source line 4227
 \begin{proof}
  \proves{thm:simple-triangle}
  \leanok
  \uses{thm:g3.1.7}
  According to~\eqref{eq:cond-addit}, we have
\[
  K(y, z \mvert x) \eqa K(y \mvert x) + K(z \mvert y, K(y \mvert x), x).
\]
The left-hand side is \( \gea K(z \mvert x) \), the second term of the
right-hand side is \( \lea K(z \mvert y) \).
 \end{proof}

Equation~\eqref{eq:cond-addit} justifies the following.
% Gacs source line 4237
 \begin{definition}
  \label{def:g3.1.14}
  \lean{Kolmogorov.Information.CondInfoStar}
  \leanok
  \uses{def:K, def:pair-complexity}
 \begin{equation}\label{eq:cond-inf}
 I^{*}(x : y \mvert z) = K(x \mvert z) + K(y \mvert z) - K(x, y \mvert z).
 \end{equation}
 \end{definition}

Then we have
 \begin{align*}
    I^{*}(x : y \mvert z) &\eqa K(y \mvert z) - K(y \mvert x, K(x \mvert z), z)
\\                      &\eqa K(x \mvert z) - K(x \mvert y, K(y \mvert z), z).
 \end{align*}

\begin{formalization}
\texttt{Kolmogorov.Information.condInfoStar\_eq\_left} and
\texttt{Kolmogorov.Information.condInfoStar\_eq\_right}, from Theorem~\ref{thm:g3.1.7}.
\end{formalization}

% Gacs source line 4249
 \begin{theorem}[Algorithmic data processing identity]
  \label{thm:g3.1.9}
  \lean{Kolmogorov.Information.data_processing}
  \leanok
  \uses{def:I*, def:g3.1.14}
 We have
 \begin{equation}\label{eq:alg-data-proc-ident}
   I^{*}(z : (x, y)) \eqa I^{*}(z : x) + I^{*}(z : y \mvert x^{*}).
 \end{equation}
 \end{theorem}

% Gacs source line 4255
 \begin{proof}
  \proves{thm:g3.1.9}
  \leanok
  \uses{thm:g3.1.4}
We have
 \begin{align*}
    I^{*}(z : (x, y)) &\eqa K(x, y) + K(z) - K(x, y, z)
\\  I^{*}(z : x)      &\eqa K(x)    + K(z) - K(x, z)
\\ I^{*}(z : (x, y)) - I^{*}(z : x) 
            &\eqa K(x, y) + K(x, z) - K(x, y, z) - K(x)
\\          &\eqa K(y \mvert x^{*}) + K(z \mvert x^{*}) - K(y, z \mvert x^{*})
\\          &\eqa I^{*}(z : y \mvert x^{*}),
 \end{align*}
where we used additivity and the definition~\eqref{eq:cond-inf} repeatedly.
 \end{proof}

\begin{formalization}
The additivity used is \eqref{eq:Haddit.alt}, three times; the Lean proof needs no
conditional addition.
\end{formalization}

% Gacs source line 4267
 \begin{corollary}
  \label{cor:g3.1.15}
  \lean{Kolmogorov.Information.infoStar_mono}
  \leanok
  \uses{def:I*}
The information \( I^{*} \) is monotonic:
 \begin{equation}\label{eq:inf-mon}
   I^{*}(x : (y, z)) \gea I^{*}(x : y).
 \end{equation}
  \end{corollary}

% Proof environment ours: Gacs states the corollary without proof.
\begin{proof}
  \proves{cor:g3.1.15}
  \leanok
  \uses{thm:g3.1.9}
By Theorem~\ref{thm:g3.1.9}, \( I^{*}(x : (y, z)) \eqa I^{*}(x : y) + I^{*}(x : z \mvert y^{*}) \),
and \( I^{*}(x : z \mvert y^{*}) \gea 0 \) because
\( K(x, z \mvert y^{*}) \lea K(x \mvert y^{*}) + K(z \mvert y^{*}) \).
\end{proof}

The following theorem is also sometimes useful.

% Gacs source line 4276
 \begin{theorem}[Starry triangle inequality]\label{lem:triangle}
  \lean{Kolmogorov.Information.starry_triangle}
  \leanok
  \uses{def:K, def:pair-complexity}
For all \( x,y,z \), we have
 \begin{equation}\label{eq:starry-triangle}
  K(x \mvert y^*) \lea K(x, z \mvert y^{*})
                \lea K(z \mvert y^*) + K(x \mvert z^*).
 \end{equation}
 \end{theorem}

% Gacs source line 4283
 \begin{proof}
  \proves{lem:triangle}
  \leanok
  \uses{thm:g3.1.9, thm:g3.1.4}
Using the algorithmic data-processing
identity~\eqref{eq:alg-data-proc-ident} and the definitions, we have
 \begin{align*}
 K(z) & - K(z \mvert y^{*}) + 
 K(x \mvert z^{*}) - K(x \mvert y, K(y \mvert z^{*}), z^{*})  
\\ &\eqa I^{*}(y : z) + I^{*}(y : x \mvert z^{*}) \eqa I^{*}(y : (x, z)) 
\\ &\eqa K(x, z) - K(x, z \mvert y^{*}),
 \end{align*}
which gives, after cancelling \( K(z) + K(x \mvert z^{*}) \) on the
left-hand side with \( K(x, z) \) on the right-hand side, to which it is equal
according to additivity, and changing sign:
 \[
  K(x, z \mvert y^{*}) \eqa
  K(z \mvert y^{*}) + K(x \mvert y, K(y \mvert z^{*}), z^{*}).
 \]
From here, we proceed as in the proof of the simple triangle
inequality~\eqref{eq:simple-triangle}
\end{proof}


\subsection{Information non-increase}

In this subsection, we somewhat follow the exposition  
of~\cite{GacsTrompVitAlgorStat00}.

The classical data processing identity has a simple but important
application.
Suppose that the random variables \( Z, X, Y \) form a Markov chain in this
order.
Then \( I(Z : Y \mvert X) = 0 \) and hence \( I(Z : (X, Y)) = I(Z : X) \).
In words: all information in \( Z \) about \( Y \) is coming through \( X \): a Markov
transition from \( X \) to \( Y \) cannot increase the information about \( Z \).
Let us try to find the algorithmic counterpart of this theorem.

Here, we rigorously show that this is the case in the algorithmic statistics
setting:
the information in one object about another
cannot be increased by any deterministic algorithmic method
by more than a constant. 
With added randomization this holds with overwhelming probability.
For more elaborate versions see~\cite{Levin74,LevinRandCons84}.

Suppose we want to obtain information about a certain object \( x \). 
It does not seem to be a good policy to guess blindly. 
This is confirmed by the following inequality.
 \begin{equation}\label{eq:info-guess}
	\sum_y \m(y)2^{I(x : y)} \lem 1 
 \end{equation}
  which says that for each \( x \), the expected value of \( 2^{I(x : y)} \) is 
small, even with respect to the universal constructive semimeasure 
\( \m(y) \). 
The proof is immediate if we recognize that by the Coding 
Theorem, 
 \[	2^{I(x : y)} \eqm \frac{2^{-K(y \mvert x)}}{\m(y)}
 \]
  and use the estimate~\eqref{eq:condHsum}.

\begin{formalization}
\eqref{eq:info-guess} is \texttt{Kolmogorov.Information.info\_guess}, for any universal
semimeasure \( \m \). With \( 2^{-K(y)} \) in place of \( \m(y) \) the sum equals
\( \sum_{y} 2^{-K(y \mvert x)} \le 1 \) exactly
(\texttt{Kolmogorov.Information.info\_guess\_weight}).
\end{formalization}

We prove a strong version of the information non-increase law
under deterministic processing (later we need the attached corollary):

% Gacs source line 4344
\begin{theorem}
  \label{thm:g3.1.11}
  \lean{Kolmogorov.Information.info_nonincrease_det}
  \leanok
  \uses{def:I*}
Given \( x \) and \( z \), let \( q \) be a program computing \( z \) from \( x^{*} \).
Then we have
 \begin{equation*}
   I^{*}(y : z) \lea I^{*}(y : x) + K(q).
 \end{equation*}
\end{theorem}

\begin{formalization}
Programs are those of an arbitrary fixed partial recursive programming system
\( \Phi(q, x^{*}) = z \), and the constant depends only on \( \Phi \).
\end{formalization}

% Gacs source line 4351
\begin{proof}
  \proves{thm:g3.1.11}
  \leanok
  \uses{cor:g3.1.15, thm:g3.1.9}
By monotonicity~\eqref{eq:inf-mon} and
the data processing identity~\eqref{eq:alg-data-proc-ident}:
\[
 I^{*}(y : z) \lea
I^{*}(y : (x, z)) \eqa I^{*}(y : x) + I^{*}(y : z \mvert x^{*}).   
\]
By definition of information and the definition of \( q \):
 \[
   I^{*}(y : z \mvert x^{*}) \lea K(z \mvert x^{*}) \lea K(q).
 \]
\end{proof}

Randomized computation can increase information only with negligible
probability.
Suppose that \( z \) is obtained from \( x \) by some randomized computation.
The probability \( p(z \mvert x) \) of obtaining \( z \) from \( x \) is a semicomputable
distribution over the \( z \)'s.
The information increase \( I^{*}(z : y) - I^{*}(x : y) \) satisfies the
theorem below. 

% Gacs source line 4372
 \begin{theorem}
  \label{thm:g3.1.12}
  \lean{Kolmogorov.Information.info_nonincrease_rand_mul, Kolmogorov.Information.info_nonincrease_rand}
  \leanok
  \uses{def:I*, def:m}
For all \( x, y, z \) we have
 \[
  \m(z \mvert x^{*}) 2^{I^{*}(z : y) - I^{*}(x : y)}
  \lem \m(z \mvert x^{*}, y, K(y \mvert x^{*})).
 \]
Therefore it is upperbounded by \( \m(z \mvert x) \lem \m(z \mvert x^{*}) \).
 \end{theorem}

\begin{formalization}
The conditional universal semimeasure is taken in its coding form
\( \m(z \mvert w) = 2^{-K(z \mvert w)} \), equal to it within a constant factor by the
conditional coding theorem; the condition \( x^{*}, y, K(y \mvert x^{*}) \) is the string
\( \ang{x^{*}, y, K(y \mvert x^{*})} \). The additive form,
\( I^{*}(z : y) - I^{*}(x : y) \lea K(z \mvert x^{*}) - K(z \mvert x^{*}, y, K(y \mvert x^{*})) \),
is \texttt{Kolmogorov.Information.info\_nonincrease\_rand}. The sentence after the display is
not formalized.
\end{formalization}

% Gacs source line 4380
\begin{remark}
  \label{rem:g3.1.16}
The theorem implies
 \begin{equation}\label{eq:inf-incre-exp}
   \sum_{z} \m(z \mvert x^{*}) 2^{I^{*}(z : y) - I^{*}(x : y)} \lem 1.
 \end{equation}
This says that
the \( \m(\cdot \mvert x^{*}) \)-expectation of the exponential of the increase
is bounded by a constant.
It follows that for example, 
the probability of an increase of mutual information
by the amount \( d \) is \( \lem 2^{-d} \).
\end{remark}

\begin{formalization}
With \( \m(z \mvert x^{*}) = 2^{-K(z \mvert x^{*})} \): \eqref{eq:inf-incre-exp} is
\texttt{Kolmogorov.Information.info\_increase\_expectation}, and the bound \( \lem 2^{-d} \)
on the weight of the \( z \) that increase the information by at least \( d \) is
\texttt{Kolmogorov.Information.info\_increase\_tail}.
\end{formalization}

% Gacs source line 4393
 \begin{proof}[Proof of Theorem~\protect\ref{thm:g3.1.12}]
  \proves{thm:g3.1.12}
  \leanok
  \uses{thm:g3.1.9, cor:g3.1.15}
We have, by the data processing inequality:
 \[
 I^{*}(y : (x, z)) - I^{*}(y : x) 
   \eqa I^{*}(y : z \mvert x^{*}) 
   \eqa K(z \mvert x^{*}) - K(z \mvert y, K(y \mvert x^{*}), x^{*}).
 \]
Hence, using also \( I^{*}(y: (x, z)) \gea I^{*}(y : z) \) by monotonicity:
 \[
 I^{*}(y : z) - I^{*}(y : x) - K(z \mvert x^{*}) 
   \lea - K(z \mvert y, K(y \mvert x^{*}), x^{*}).
 \]
Putting both sides into the exponent gives the statement of the theorem.
 \end{proof}

% Gacs source line 4408
\begin{remark}
  \label{rem:g3.1.17}
The theorem on information non-increase and its proof look similar to
the theorems on randomness conservation.
There is a formal connection, via the observation made in
Remark~\ref{rem:inf-def-rand}.
Due to the difficulties mentioned there, we will explore the connection
only later in our notes.
\end{remark}


\section{The complexity of decidable and enumerable sets}\label{sec:halting}

If \( \og = \og(1)\og(2)\dotsb \) is an infinite binary sequence, then we may
be interested in how the complexity of the initial segments
 \[
  \og(\frto{1}{n}) = \og(1) \dotsb \og(n)
 \]
grows with \( n \).
We would guess that if \( \og \) is ``random'' then \( C(\og(\frto{1}{n})) \) is
approximately \( n \); this question will be studied satisfactorily in later
sections.
Here, we want to look at some other questions.

Can we determine whether the sequence \( \og \) is computable, by just looking
at the complexity of its initial segments?
It is easy to see that if \( \og \) is computable then 
\( C(\og(\frto{1}{n})) \lea \log n \).
But of course, if \( C(\og(\frto{1}{n})) \lea \log n \) then it is not necessarily true
yet that \( \og \) is computable.
Maybe, we should put \( n \) into the condition.
If \( \og \) is computable then \( C(\og(\frto{1}{n}) \mvert n) \lea 0 \).
Is it true that \( C(\og(\frto{1}{n}) \mvert n) \lea 0 \) for all \( n \) then indeed
\( \og \) is computable?
Yes, but the proof is quite difficult (see either its original, attributed
to Albert Meyer in~\cite{Loveland69}, or in~\cite{ZvLe70}).
The suspicion arises that when
measuring the complexity of starting segments of infinite sequences,
neither Kolmogorov complexity nor its prefix version are the most natural
choice.
Other examples will support the suspicion, so let us introduce, following
Loveland's paper~\cite{Loveland69}, the following variant. 

% Gacs source line 4454
 \begin{definition}
  \label{def:g3.2.1}
  \lean{Kolmogorov.Information.Halting.DecMap, Kolmogorov.Information.Halting.Decides, Kolmogorov.Information.Halting.decK}
  \leanok
Let us say that a program \( p \) \df{decides} a string \( x = (x(1),x(2),\dotsc) \)
on a Turing machine \( T \) up to \( n \) if for all \( i \in \{1,\dots,n\} \) we have
 \[
   T(p,i) = x(i).
 \]
Let us further define the \df{decision complexity}
 \[
   C_{T}(x; n) = \min\setof{\len{p} : p \txt{ decides \( x \) on \( T \) up to \( n \)}}.
 \]
 \end{definition}

\begin{formalization}
A machine is a partial recursive map \( T \colon \dS \times \dN \to \{0, 1\} \), and positions are
counted from \( 0 \). The decision complexity is defined for \( n = \len{x} \) only:
\texttt{decK T x} is \( C_{T}(x; \len{x}) \), which is how the notation is used below.
\end{formalization}

As for the Kolmogorov complexity, an invariance theorem holds, and there is
an optimal machine \( T_{0} \) for this complexity.
Again, we omit the index \( T \), assuming that such an optimal machine has
been fixed.

The differences between \( C(x) \), \( C(x \mvert n) \) and \( C(x ; n) \) are
interesting to explore; intuitively, the important difference is that the
program achieving the decision complexity of \( x \) does not have to offer
precise information about the length of \( x \).
We clearly have
 \[
   C(x \mvert n) \lea C(x ; n) \lea C(x).
 \]
Examples that each of these inequalities can be strict, are left to
exercises.

\begin{formalization}
Optimality is \texttt{Kolmogorov.Information.Halting.IsOptimalDecision}, and the invariance
theorem is \texttt{Kolmogorov.Information.Halting.univDec\_isOptimal}: the universal
decompressor, run on condition \( i \), with its first output bit kept, is an optimal decision
machine. The two inequalities are
\texttt{Kolmogorov.Information.Halting.condK\_le\_decK} (with \( n = \len{x} \) given in
unary) and \texttt{Kolmogorov.Information.Halting.decK\_le\_plainK}.
\end{formalization}

% Gacs source line 4483
 \begin{remark}
  \label{rem:g3.2.2}
If \( \og(\frto{1}{n}) \) is a segment of an infinite sequence \( \og \) then we can write
 \[
   C(\og ; n)
 \]
instead of \( C(\og(\frto{1}{n}) ; n) \) without confusion, since there is only one way
to understand this expression.
 \end{remark}

\begin{formalization}
An infinite sequence is a map \( \og \colon \dN \to \{0, 1\} \), and \( \og(\frto{1}{n}) \) is
\texttt{Kolmogorov.Information.Halting.seqPrefix}.
\end{formalization}

Decision complexity offers an easy characterization of decidable infinite
sequences.

% Gacs source line 4495
 \begin{theorem}
  \label{thm:g3.2.1}
  \lean{Kolmogorov.Information.Halting.computable_iff_decK_bounded}
  \leanok
  \uses{def:g3.2.1}
Let \( \og = (\og(1), \og(2),\dotsc) \) be an infinite sequence.
Then \( \og \) is decidable if and only if
\begin{equation}\label{eq:dec-compl-0}
   C(\og ; n) \lea 0.
\end{equation}
 \end{theorem}

\begin{formalization}
``Decidable'' is Mathlib's \texttt{Computable}.
\end{formalization}

% Gacs source line 4503
 \begin{proof}
  \proves{thm:g3.2.1}
  \leanok
If \( \og \) is decidable then~\eqref{eq:dec-compl-0} is obviously true.
Suppose now that~\eqref{eq:dec-compl-0} holds:
there is a constant \( c \) such that for all \( n \) we have \( C(\og(\frto{1}{n});n) < c \).
Then there is a program \( p \) of length \( \le c \) with the property that for
infinitely many \( n \), decides \( x \) on the optimal machine \( T_{0} \) up to \( n \).
Then for all \( i \), we have \( T_{0}(p,i) = \og(i) \), showing that \( \og \) is
decidable.
 \end{proof}

Let us use now the new tool to prove a somewhat more surprising result.
Let \( \og \) be a 0-1 sequence that is the indicator function of a recursively
enumerable set.
In other words, the set \( \setof{i: \og(i) = 1} \) is recursively enumerable.
As we know such sequences are not always decidable, so \( C(\og(\frto{1}{n}) ; n) \) is
not going to be bounded.
How fast can it grow?
The following theorem gives exact answer, showing that it grows
only logarithmically.

% Gacs source line 4524
 \begin{theorem}[Barzdin, see~\cite{Barzdin68}]
  \label{thm:g3.2.2}
  \lean{Kolmogorov.Information.Halting.decK_le_log_of_isRE, Kolmogorov.Information.Halting.exists_isRE_log_le_decK}
  \leanok
  \uses{def:g3.2.1}
 \begin{alphenum}
  \item\label{i:Barzdin-upper}
Let \( \og \) be the indicator function of a recursively enumerable set \( E \).
Then we have
 \[
   C(\og ; n) \lea \log n.
 \]
  \item\label{i:Barzdin-lower}
The set \( E \) can be chosen such that for all \( n \) we have 
\( C(\og ; n) \ge \log n \).
 \end{alphenum}
 \end{theorem}

\begin{formalization}
Logarithms are \( \lfloor \log_{2} \cdot \rfloor \). Part~(b) holds against every partial
recursive decision machine, optimal or not.
\end{formalization}

% Gacs source line 4537
 \begin{proof}
  \proves{thm:g3.2.2}
  \leanok
Let us prove~\eqref{i:Barzdin-upper} first.
Let \( n' \) be the first power of 2 larger than \( n \).
Let \( k(n) = \sum_{i \le n'} \og(i) \).
For a constant \( d \), let \( p = p(n, d) \) be a program
that contains an initial segment \( q \) of size \( d \) followed by a
the number \( k(n) \) in binary notation, and padded to \( \cei{\log n} + d \) 
with 0's.
The program \( q \) is self-delimiting, so it sees where \( k(n) \) begins.

The machine \( T_{0}(p, i) \) works as follows, under the control of \( q \).
From the length of the whole \( p \), it determines \( n' \).
It begins to enumerate \( \og(\frto{1}{n}') \) until it found \( k(n) \) 1's.
Then it knows the whole \( \og(\frto{1}{n}') \), so it outputs \( \og(i) \).

Let us prove now~\eqref{i:Barzdin-lower}.
Let us list all possible programs \( q \) for the machine \( T_{0} \) as 
\( q = 1, 2, \dotsc \).
Let \( \og(q) = 1 \) if and only if \( T_{0}(q, q) = 0 \).
The sequence \( \og \) is obviously the indicator function of a recursively
enumerable set.
To show that \( \og \) has the desired property, assume that for some
\( n \) there is a \( p \) with \( T_{0}(p, i) = \og(i) \) for all \( i \le n \).
Then \( p > n \) since \( \og(p) \) is defined to be different from \( T_{0}(p, p) \).
It follows that \( \len{p} \ge \log n \).
 \end{proof}

\begin{formalization}
In~\eqref{i:Barzdin-upper} the program is the number \( k \) of elements of \( E \) below
\( N = 2^{L} - 1 \), written in \( L = \lfloor \log n \rfloor + 1 \) bits; the machine reads
\( N \) off the length of the program, so no prefix \( q \) is needed. In~\eqref{i:Barzdin-lower}
a program \( p \) is identified with a number in \( [2^{\len{p}}, 2^{\len{p}+1}) \)
(\texttt{Kolmogorov.Information.Halting.idx}), one-to-one.
\end{formalization}

Decision complexity has been especially convenient for the above theorem.
Neither \( C(\og(\frto{1}{n})) \) nor \( C(\og(\frto{1}{n}) \mvert n) \) would have been suitable to
formulate such a sharp result.
To analyze the phenomenon further, we introduce some more concepts.

% Gacs source line 4569
 \begin{definition}
  \label{def:g3.2.3}
  \lean{Kolmogorov.Information.Halting.haltSet, Kolmogorov.Information.Halting.chi, Kolmogorov.Information.Halting.haltsIn}
  \leanok
  \uses{def:K}
Let us denote, for this section, by \( E \) the set of those binary strings \( p \) on
which the optimal prefix machine \( T \) halts:
 \begin{equation}\label{eq:E^{t}}
 \begin{split}
     E^{t} &= \setof{p : T(p) \txt{ halts in \( < t \) steps} },
\\   E     &= E^{\infty}.
 \end{split}
 \end{equation}
Let 
 \begin{equation}\label{eq:chi}
  \chi = \chi_{E}
 \end{equation}
be the infinite sequence that is the indicator function of the
set \( E \), when the latter is viewed as a set of numbers.
 \end{definition}

\begin{formalization}
\( E \) is the domain of the prefix machine on the empty condition. ``Halts in \( < t \) steps''
is read as ``halts within fuel \( t \)'' for a fixed code of the machine
(\texttt{Kolmogorov.Information.Halting.haltsIn}, Mathlib's \texttt{Nat.Partrec.Code.evaln}),
and \( E^{t} \) contains only programs of length at most \( t \). A string \( p \) is viewed as
the number \( 2^{\len{p}} + \sum_{i} p_{i} 2^{i} \) (\texttt{Kolmogorov.Information.Halting.idx}),
so the strings of length \( k \) get the numbers in \( [2^{k}, 2^{k+1}) \).
\end{formalization}

It is easy to see that the set \( E \) is complete among recursively enumerable
sets with respect to many-one reduction.
The above theorems show that though it contains an infinite amount of
information, this information is not stored in the sequence \( \chi \) densely at
all: there are at most \( \log n \) bits of it in the segment \( \chi(\frto{1}{n}) \).
There is an infinite sequence, though, in which the same information is
stored much more densely: as we will see later, maximally densely.

% Gacs source line 4594
 \begin{definition}[Chaitin's Omega]
  \label{def:g3.2.4}
  \lean{Kolmogorov.Information.Halting.omegaStage, Kolmogorov.Information.Halting.omega}
  \leanok
  \uses{def:g3.2.3}
Let
 \begin{equation}\label{eq:Og}
 \begin{split}
   \Og^{t} &= \sum_{p \in E^{t}} 2^{-\len{p}},
\\ \Og     &= \Og^{\infty}.
 \end{split}
 \end{equation}
 \end{definition}

\begin{formalization}
\( \Og \) is the halting probability of the prefix machine on the empty condition, a value in
\( [0, 1] \), and \( \Og = \sup_{t} \Og^{t} \)
(\texttt{Kolmogorov.Information.Halting.omega\_eq\_iSup}).
\end{formalization}

Let \( \Og(\frto{1}{n}) \) be the sequence of the first \( n \) binary digits in the
expansion of \( \Og \), and let it also denote the binary number
\( 0.\Og(1)\dotsb\Og(n) \).
Then we have
 \[
    \Og(\frto{1}{n}) < \Og < \Og(\frto{1}{n}) + 2^{-n}.
 \]

\begin{formalization}
\( \Og(\frto{1}{n}) \) is \texttt{Kolmogorov.Information.Halting.omegaPrefix}, the \( n \)-bit
binary notation of \( a_{n} = \min(\lfloor \Og 2^{n} \rfloor, 2^{n} - 1) \). The bounds are
proved with \( \le \) (\texttt{Kolmogorov.Information.Halting.omegaDigits\_le},
\texttt{Kolmogorov.Information.Halting.omega\_le\_omegaDigits\_succ}). Strict inequality
needs \( \Og \) not to be a dyadic rational. That holds for an exactly universal machine
(see Theorem~\ref{thm:g3.2.3}), and is proved as
\texttt{Kolmogorov.Information.Turing.omega\_ne\_dyadic}.
\end{formalization}

% Gacs source line 4612
 \begin{theorem}
  \label{thm:g3.2.3}
  \lean{Kolmogorov.Information.Turing.chi_turingEquivalent_omega, Kolmogorov.Information.Turing.chi_le_omega, Kolmogorov.Information.Turing.omega_le_chi}
  \leanok
  \uses{def:g3.2.3, def:g3.2.4}
The sequences \( \Og \) and \( \chi \) are Turing-equivalent.
 \end{theorem}

\begin{formalization}
The optimal machine \( T \) is assumed \emph{exactly universal}: every prefix machine
\( M \) is simulated as \( T(tp) = M(p) \) for a fixed tag \( t \)
(\texttt{Kolmogorov.Information.Oracle.IsExactUniversal}). The repository's universal
machine has this property
(\texttt{Kolmogorov.Information.Oracle.universalPrefixMachine\_isExactUniversal}). Turing
equivalence is Mathlib's \texttt{TuringEquivalent} of \( \chi \) and of the digit sequence of
\( \Og \) (\texttt{Kolmogorov.Information.Turing.omegaBit}, digits counted from \( 0 \)),
both as partial functions \( \dN \to \dN \).
\end{formalization}

% Gacs source line 4615
 \begin{proof}
  \proves{thm:g3.2.3}
  \leanok
Let us show first that given \( \Og \) as an oracle, a Turing machine can
compute \( \chi \).
Suppose that we want to know for a given string \( p \) of length \( k \) 
whether \( \chi(p) = 1 \) that is whether \( T(p) \) halts.
Let \( t(n) \) be the first \( t \) for which \( \Og^{t} > \Og(\frto{1}{n}) \).
If a program \( p \) of length \( n \) is not in \( E^{t(n)} \) then it is not in \( E \)
at all, since \( 2^{-\len{p}} = 2^{n} < \Og - \Og^{t} \).
It follows that \( \chi(1 : 2^{n}) \) can be computed from \( \Og(\frto{1}{n}) \).

To show that \( \Og \) can be computed from \( E \), let
us define the recursively enumerable set
 \[
   E' = \setof{r : r \txt{ rational, } \Og > r}.
 \]
The set \( E' \) is reducible to \( E \) since the latter
is complete among recursively enumerable sets
with respect to many-one reduction.
On the other hand, \( \Og \) is obviously computable from \( E' \).
 \end{proof}

\begin{formalization}
For \( \chi \le_{T} \Og \) the Lean proof takes \( t \) to be the first stage with
\( \Og^{t} + 2^{-n} > a_{n} 2^{-n} \), which needs no irrationality of \( \Og \); this direction
holds for every prefix machine. For \( \Og \le_{T} \chi \), many-one completeness of \( E \)
comes from exact universality. ``\( \Og \) is computable from \( E' \)'' uses
\( \lfloor \Og 2^{N} \rfloor = \card{\setof{k < 2^{N} : 0 < k,\ k 2^{-N} < \Og}} \), which holds
because \( \Og \) is not dyadic. That in turn follows because \( \chi \le_{T} \Og \) and \( \chi \)
is not computable (\texttt{Kolmogorov.Information.Turing.chi\_not\_computable}). Reductions are
built from a use principle for Mathlib's \texttt{RecursiveIn}
(\texttt{Kolmogorov.Information.Turing.turingReducible\_of\_table}).
\end{formalization}

The following theorem shows that \( \Og \) stores the information more densely.

% Gacs source line 4638
 \begin{theorem}\label{thm:Og-complex}
  \lean{Kolmogorov.Information.Halting.le_KP_omegaPrefix}
  \leanok
  \uses{def:g3.2.4, def:K}
 We have \( K(\Og(\frto{1}{n})) \gea n \).
 \end{theorem}

% Gacs source line 4641
 \begin{proof}
  \proves{thm:Og-complex}
  \leanok
Let \( p_{1} \) be a self-delimiting program outputting \( \Og(\frto{1}{n}) \).
Recall the definition of the sets \( E^{t} \) in~\eqref{eq:E^{t}}
and the numbers \( \Og^{t} \) in~\eqref{eq:Og}.
Let \( \Og_{1} \) denote the real number whose binary digits after 0 are given
by \( \Og(\frto{1}{n}) \), and let \( t_{1} \) be the first \( t \) with \( \Og^{t}>\Og_{1} \).
Let \( x_{1} \) be the first string \( x \) such that \( T(p)\ne x \) for any
\( p\in E^{t_{1}} \).

We have \( K(x_{1})\ge n \).
On the other hand, we just computed \( x_{1} \) from \( p_{1} \), so
\( K(x_{1}\mvert p_{1})\lea 0 \).
We found \( n \le K(x_{1}) \lea K(p_{1}) + K(x_{1}\mvert p_{1}) \lea K(p_{1}) \).
 \end{proof}

\begin{formalization}
Holds for every optimal prefix machine. The Lean proof takes \( t_{1} \) to be the first stage
with \( \Og^{t} + 2^{-n} > \Og_{1} \), so it needs no irrationality of \( \Og \): then every
halting program of length less than \( n \) is in \( E^{t_{1}} \)
(\texttt{Kolmogorov.Information.Halting.halts\_of\_short}). The string \( x_{1} \) is searched
among \( 0^{j} \), \( j = 0, 1, \dotsc \).
\end{formalization}


\section{The complexity of complexity}\label{sec:compl-compl}

\subsection{Complexity is sometimes complex}

This section is devoted to a quantitative estimation of the uncomputability
of the complexity function \( C(x) \).
Actually, we will work with the prefix complexity \( K(x) \), but the results
would be very similar for \( C(x) \).
The first result shows that the value
\( K(x) \) is not only not computable from \( x \), but
its conditional complexity \( K(K(x) \mvert x) \) given \( x \) is sometimes quite
large.
How large can it be expected to be?
Certainly not much larger than \( \log K(x) + 2\log\log K(x) \), since we can
describe any value \( K(x) \) using this many bits.
But it can come close to this, as shown by Theorem~\ref{thm:compl-compl}.
This theorem says that for all \( n \), there exists \( x \) of length \( n \) with
 \begin{equation}\label{eq:compl-compl}
  K(K(x) \mvert x) \gea \log n - \log \log n.
 \end{equation}

% Gacs source line 4678
 \begin{proof}[Proof of Theorem~\protect\ref{thm:compl-compl}]
  \proves{thm:compl-compl}
  \leanok
Let \( U(p, x) \) be the optimal self-delimiting machine for which
\( K(y \mvert x) = \min_{U(p, x) = y} \len{p} \).
Let \( s \le \log n \) be such that if \( \len{x} = n \) then 
a \( p \) of length \( x \) can be found for which \( U(p, x) = K(x) \).
We will show
 \begin{equation}\label{eq:compl-compl.s-lb}
   s \gea \log n - \log \log n.
 \end{equation}
Let us say that \( p \in \{0, 1\}^{s} \) is \df{suitable} for 
\( x \in \{0,1\}^{n} \) if there exists a \( k = U(p, x) \) and a
\( q \in \{0,1\}^{k} \) with \( U(p, \Lambda) = x \).
Thus, \( p \) is suitable for \( x \) if it produces the length of some program of
\( x \), not necessarily a shortest program.

Let \( M_{i} \) denote the set of those \( x \) of length \( n \)
for which there exist at least \( i \)
different suitable \( p \) of length \( s \).
We will examine the sequence
 \[
   \{0,1\}^{n} = M_{0} \spsq M_{1} \spsq \dots 
   \spsq M_{j} \spsq M_{j+1} = \emptyset,
 \]
where \( M_{j} \ne \emptyset \).
It is clear that \( 2^{s} \ge j \).
To lowerbound \( j \), we will show that the sets \( M_{i} \) decrease only slowly:
 \begin{equation}\label{eq:compl-compl.slow-decr}
  \log \card{M_{i}} \lea \log\card{M_{i+1}} + 4 \log n.
 \end{equation}
We can assume 
 \begin{equation}\label{eq:compl-compl.assume}
   \log\card{M_{i} \xcpt M_{i+1}} \ge \log\card{M_{i}} - 1,
 \end{equation}
otherwise~\eqref{eq:compl-compl.slow-decr} is trivial.
We will write a program that finds an element \( x \) of 
\( M_{i} \xcpt M_{i+1} \) with the property \( K(x) \ge \log\card{M_{i}} - 1 \).
The program works as follows.
 \begin{itemize}

  \item It finds \( i,n \) with the help of descriptions of length
\( \log n + 2\log\log n \), and \( s \) with the help of a description of length
\( 2 \log\log n \).

  \item It finds  \( \card{M_{i+1}} \) with the help of a description of length
\( \log\card{M_{i+1}} + \log n + 2\log\log n \).

  \item From these data, it determines the set \( M_{i+1} \), and then begins
to enumerate the set \( M_{i} \xcpt M_{i+1} \) as \( x_{1}, x_{2},\dotsc \).
For each of these elements \( x_{r} \), it knows that there are exactly \( i \)
programs suitable for \( x_{r} \), find all those, and find the shortest
program for \( x \) produced by these.
Therefore it can compute \( K(x_{r}) \).

  \item According to the assumption~\eqref{eq:compl-compl.assume},
there is an \( x_{r} \) with \( K(x_{r}) \ge \log\card{M_{i}} - 1 \).
The program outputs the first such \( x_{r} \).

 \end{itemize}
The construction of the program shows that its length is
\( \lea \log\card{M_{i+1}} + 4\log n \), hence for the \( x \) 
we found 
 \[
   \log\card{M_{i}} - 1 \le K(x) \lea \log\card{M_{i+1}} + 4\log n,
 \]
which proves~\eqref{eq:compl-compl.slow-decr}.
This implies \( j \ge n/(4 \log n) \), and hence~\eqref{eq:compl-compl.s-lb}.
 \end{proof}

\begin{formalization}
The Lean proof makes the sketch explicit. \( s \) is \( \max_{\len{x} = n} K(K(x) \mvert x) \),
and \( p \) with \( \len{p} \le s \) is suitable for \( x \) if \( U(p, x) \) names an upper bound on
\( K(x) \) (\texttt{Kolmogorov.Information.CompComp.Suit}). The program that finds \( x \) is a
search machine on the five numbers \( n, s, i, \card{M_{i+1}}, m \)
(\texttt{Kolmogorov.Information.CompComp.dec\_spec}). Assumption~\eqref{eq:compl-compl.assume}
is not needed: with \( m = \lfloor \log \card{M_{i} \xcpt M_{i+1}} \rfloor \) the machine outputs
some \( x \) with \( K(x) \ge m \). This gives
\( \card{M_{i}} + 1 \le (A + 1)(\card{M_{i+1}} + 1) \) with
\( A = 2^{14 \log n + O(1)} \) (\texttt{Kolmogorov.Information.CompComp.slow\_decrease}); the
constant \( 4 \) becomes \( 14 \) because the five numbers are coded with
\( K(v) \le \log v + 2 \log\log v + O(1) \), which does not affect the result. Iterating from
\( \card{M_{1}} = 2^{n} \) to \( \card{M_{2^{s+1}}} = 0 \) gives \( \log n \le s + \log\log n + O(1) \).
\end{formalization}


\subsection{Complexity is rarely complex}

Let
 \[
  f(x) = K(K(x) \mvert x).
 \]
We have seen that \( f(x) \) is sometimes large.
Here, we will show that the sequences \( x \) for which this happens are rare.
Recall that we defined \( \chi \) in~\eqref{eq:chi} as
the indicator sequence of the halting problem.

% Gacs source line 4758
 \begin{theorem}\label{thm:co-rarely-co}
  \lean{Kolmogorov.Information.Oracle.rarely_complex, Kolmogorov.Information.Oracle.rarely_complex_info}
  \leanok
  \uses{def:g3.1.8, def:g3.2.3}
We have
 \[
   I(\chi : x) \gea f(x) - 2.4 \log f(x).
 \]
 \end{theorem}

\begin{formalization}
We prove a sharper statement. There is a constant \( c \) with
 \[
   K(x \mvert \chi) + f(x) \le K(x) + 2K(f(x)) + c
   \quad\text{for all } x,
 \]
that is, \( I(\chi : x) \ge f(x) - 2K(f(x)) - c \), where \( K(f(x)) \) is the
complexity of the number \( f(x) \). This implies G\'acs' version. Since
\( K(n) \le \log n + 2\log\log n + O(1) \le 1.2 \log n + O(1) \), we get
\( 2K(f(x)) \le 2.4 \log f(x) + O(1) \). His constant \( 2.4 = 2 \times 1.2 \)
was chosen to make the sum in his proof converge; the sharp form needs no such choice.

Further departures from the printed statement:
 \begin{itemize}
 \item The interpreter \( U \) is assumed \emph{exactly universal}: every prefix
  interpreter \( M \) is simulated as \( U(tp, y) = M(p, y) \) for a fixed tag \( t \)
  (\texttt{Kolmogorov.Information.Oracle.IsExactUniversal}). The repository's universal
  machine has this property.
 \item \( K(x \mvert \chi) \) is complexity relative to the infinite oracle \( \chi \): the
  infimum over \( m \) of \( K_V(x \mvert \chi(1:m)) \) for an optimal oracle machine \( V \).
  An oracle machine is required to be monotone in the oracle prefix. Lean name:
  \texttt{Kolmogorov.Information.Oracle.KPrel}, applied as \texttt{KPrel V (chi U) x}.
 \item The proof reads \( K(x) \), \( f(x) \) and \( K(f(x)) \) directly off finite prefixes
  of \( \chi \), so the equivalence of \( \Og \) and \( \chi \) is not used. The
  semimeasure \( \m(x) 2^{f(x)} f(x)^{-2.4} \) is replaced by
  \( \m(x) 2^{f(x)} 2^{-2K(f(x))} \).
 \end{itemize}
\end{formalization}

In view of the inequality~\eqref{eq:info-guess}, this shows that such
sequences are rare, even in terms of the universal probability, so
they are certainly rare if we measure them by any
computable distribution.
So, we may even call such sequences ``exotic''.

In the proof, we
start by showing that the sequences in question are rare.
Then the theorem will follow when we make use of the fact that
\( f(x) \) is computable from \( \chi \).

We need a lemma about the approximation of one measure by another from
below.

% Gacs source line 4778
  \begin{lemma}\label{lem:measure-approx}
  \lean{Kolmogorov.Information.Halting.measure_approx_le, Kolmogorov.Information.Halting.measure_approx_lt}
  \leanok
For any semimeasure \( \nu \) and measure \( \mu \le \nu \), let
\( S_{m} = \setof{x : 2^{m} \mu(x) \le \nu(x)} \).
Then \( \mu(S_{m}) \le 2^{-m} \).
If \( \sum_{x} (\nu(x) - \mu(x)) < 2^{-n} \) then \( \nu(S_{1}) < 2^{-n+1} \).
  \end{lemma}

\begin{formalization}
The first claim needs neither \( \mu \le \nu \) nor that \( \mu \) is a measure. In the
second, \( 2^{-n} \) is generalized to any bound \( \varepsilon \), giving
\( \nu(S_1) < 2\varepsilon \).
\end{formalization}

The proof is straightforward.

% Proof environment ours: Gacs gives none.
\begin{proof}
  \proves{lem:measure-approx}
  \leanok
\end{proof}

Let
 \[
     X(n) = \setof{x : f(x) = n}.
 \]

% Gacs source line 4793
 \begin{lemma}\label{lem:co-co-rare}
  \lean{Kolmogorov.Information.Halting.aprioriMeasure_complexityOfComplexity_le, Kolmogorov.Information.Halting.KPPlain_add_le_of_length}
  \leanok
  \uses{def:m, def:K}
 We have
 \[
  \m(X(n)) \lem n^{1.2} 2^{-n}.
 \]
 \end{lemma}

\begin{formalization}
We prove the sharper bound
 \[
   \m(X(n)) \le C\, 2^{K(n)} 2^{-n},
 \]
which implies G\'acs' version because \( 2^{K(n)} \lem n^{1.2} \). The proof below gives
the factor \( n^{1.2} \) only because it bounds
\( K(\Og(1:k)) \lea k + 1.2 \log k \). We bound the complexity of \( \Og(1:k) \) for a
well-chosen \( k \) instead. Take \( k = n - K(n) - 3e \) and describe \( \Og(1:k) \) by a
shortest program for \( n \), followed by \( e \) and by the \( k \) digits written out. Given
\( n \), \( K(n) \) and \( e \), the reader knows \( k \), so this description has length
\( n - e + O(1) \), with the \( O(1) \) independent of \( e \)
(\texttt{Kolmogorov.Information.Halting.KPPlain\_add\_le\_of\_length}). For \( x \notin Y(k) \)
this gives
\( f(x) \le n - e + c' \). Fixing \( e > c' \) yields
\( X(n) \sbsq Y(n - K(n) - 3e) \), and hence the bound.

Two smaller changes. The stage \( t(k) \) is the first \( t \) with
\( \Og^{t} + 2^{-k} > \Og(1:k) \), which makes no use of \( \Og \) being irrational. The bad
set is \( Y(k) = \setof{x : 2\m^{t(k)}(x) \le \m(x)} \), the multiplicative form of
\( K(x) \le K^{t(k)}(x) - 1 \). It satisfies \( \m(Y(k)) < 4 \cdot 2^{-k} \).
\end{formalization}

% Gacs source line 4798
 \begin{proof}
  \proves{lem:co-co-rare}
  \leanok
  \uses{lem:measure-approx}
Using the definition of \( E^{t} \) in~\eqref{eq:E^{t}}, let
\begin{align*}
      \m^{t}(x) &= \sum\setof{p : T(p) = x \txt{ in \( < t \) steps} },
\\    K^{t}(x)  &= -\log \m^{t}(x).
\end{align*}
Then according to the definition~\eqref{eq:m(x)} we have
\( \m(x) = \m^{\infty}(x) \), and \( K(x) \eqa K^{\infty}(x) \).
Let
 \begin{align*}
     t(k) &= \min\setof{t : \Og(1 : k) < \Og^{t}},
\\ \mu    &= \m^{t(k)}.
 \end{align*}
Clearly, \( K^{t(k)}(x) \gea K(x) \).
Let us show that \( K^{t(k)}(x) \)
is a good approximation for \( K(x) \) for most \( x \).
Let
 \[
   Y(k) = \setof{x : K(x) \le K^{t(k)}(x) - 1}.
 \]
By definition, for \( x \notin Y(k) \) we have
 \[
   |K^{t(k)}(x) - K(x)| \eqa 0.
 \]
On the other hand, applying Lemma~\ref{lem:measure-approx} with
\( \mu = \m^{t(k)} \), \( \nu = \m \), we obtain
 \begin{equation}\label{eq:Yk-ub}
   \m(Y(k)) < 2^{-k + 1}.
 \end{equation}
Note that
 \[
   K(K^{t(k)}(x) \mvert x, \Og(1 : k)) \eqa 0,
 \]
therefore for \( x \notin Y(k) \) we have
 \begin{align*}
  K(K(x) \mvert x, \Og(1 : k)) &\eqa 0,
\\   K(K(x) \mvert x) &\lea k + 1.2 \log k.
 \end{align*}
If \( n = k + 1.2\log k \) then \( k \eqa n - 1.2\log n \), and hence,
if \( x \notin Y(n - 1.2\log n) \) then \( K(K(x) \mvert x) \lea n \).
Thus, there is a constant \( c \) such that
 \[
   X(n) \sbsq Y(n - 1.2\log n - c).
 \]
Using~\eqref{eq:Yk-ub} this gives the statement of the lemma.
 \end{proof}

% Gacs source line 4845
 \begin{proof}[Proof of Theorem~\protect\ref{thm:co-rarely-co}]
  \proves{thm:co-rarely-co}
  \leanok
  \uses{lem:co-co-rare}
Since \( f(x) \) is computable from \( \Og \), the function
 \[
   \nu(x) = \m(x)2^{f(x)}(f(x))^{-2.4}
 \]
is computable from \( \Og \).
Let us show that it is a semimeasure (within a multiplicative constant).
Indeed, using the above lemma:
 \[
   \sum_{x} \nu(x) = \sum_{k} \sum_{x \in X(k)} \nu(x)
                 = \sum_{k} 2^{k} k^{-2.4} \m(X(k))
                 = \sum_{k} k^{-1.2} \lem 1.
 \]
Since \( \m(\cdot \mvert \Og) \) is the universal semimeasure relative to \( \Og \)
we find \( \m(x \mvert \Og) \gem \nu(x) \), hence
 \begin{align*}
   K(x \mvert \Og) &\lea -\log \nu(x) = K(x) - f(x) + 2.4\log f(x),
\\    I(\Og : x) &\gea f(x) - 2.4 \log f(x).
 \end{align*}
Since \( \Og \) is equivalent to \( \chi \), the proof is complete.
  \end{proof}

\begin{formalization}
With the sharp Lemma~\ref{lem:co-co-rare}, the same computation for
\( \nu(x) = \m(x) 2^{f(x)} 2^{-2K(f(x))} \) reads
\( \sum_x \nu(x) \le \sum_k 2^{k - 2K(k)} \m(X(k)) \le C \sum_k 2^{-K(k)} \le C \),
by Kraft's inequality. Instead of the universal semimeasure relative to \( \Og \), the Lean
proof issues the requests \( (x,\, K(x) - f(x) + 2K(f(x)) + c_0) \) to a relativized
Kraft--Chaitin coder (\texttt{Kolmogorov.Information.Oracle.exists\_oracleMachine\_of\_requests}).
The coder reads its requests off finite prefixes of \( \chi \).
\end{formalization}
